\section{Two reversors and one closure polynomial}

Put $f(q)=1-6q^2$ and
\[
H(q,p)=(f(q)-p,q),\qquad R(q,p)=(p,q),\qquad J=RH.
\]
Then $R^2=J^2=1$ and $RHR=H^{-1}$.  Parameterize $\Fix(J)$ by
\begin{equation}\label{eq:start}
q_0=X,\qquad q_{-1}=q_1=\frac{1-6X^2}{2},
\end{equation}
and continue with
\begin{equation}\label{eq:recurrence}
q_{j+1}=1-6q_j^2-q_{j-1}.
\end{equation}

For odd $n=2m+1$, define
\begin{equation}\label{eq:closure}
F_n(X)=q_{m+1}(X)-q_m(X)\in\mathbb Q[X].
\end{equation}

\begin{proposition}[Mixed-axis interpretation]\label{prop:mixed}
A root of $F_n$ gives a point in $\Fix(J)$ whose $(m+1)$st iterate lies in
$\Fix(R)$, and therefore a periodic point of period dividing $n$.
Conversely, an $n$-periodic point on $\Fix(J)$ satisfies \eqref{eq:closure}.
\end{proposition}

\begin{proof}
Let $z=(q_0,q_{-1})$.  Equation \eqref{eq:start} is exactly $Jz=z$, while
$F_n=0$ says $H^{m+1}z\in\Fix(R)$.  Since $RH^k=H^{-k}R$ and $Rz=Hz$,
the latter condition gives $H^{-m}z=H^{m+1}z$, hence $H^nz=z$.  Reversing
the argument proves the converse.
\end{proof}

This is the vertex-to-edge, or mixed-axis, closure.  It contains points of
proper odd period and may carry intersection multiplicity; both are handled
explicitly below.
